\chapter{Apr.~8 --- Quantum Groups}

\section{Intertwining Operators via Fock Modules, Continued}

\begin{theorem}
  We have (here $C = C_1 \times \cdots \times C_N$)
  \[
    \Phi^m(z)
    = \int_C \Xi(z, t_1, \dots, t_m)\,
    dt_1 \wedge \cdots \wedge dt_m.
  \]
\end{theorem}

\begin{remark}
  We also have
  \[
    \int_C \frac{\partial}{\partial t_i}
    \Xi'\, d\vec{t} = 0.
  \]
  Furethermore, $\langle \Phi(z_1) \cdots \Phi(z_N) \rangle$
  reproduce the solutions of the KZ
  equations for $|z_1| > \cdots > |z_N|$.
\end{remark}

\begin{remark}
  Let $\Phi : V_{\lambda_1, k} \to \widehat{V}_{\lambda_2, k} \otimes V_\mu(z)$.
  A \emph{Kazhdan-Lusztig (Moore-Seiberg) intertwiner} is
  \[
    \widetilde{\Phi}
    : V_{\mu, k}
    \otimes V_{\lambda_1, k}
    \longrightarrow \widehat{V}_{\lambda_2, k}(z).
  \]
  Here $V_{\mu, k}$ is a highest weight module
  for $\widehat{\g}$ and $V_\mu(z)$ is a
  lowest weight evaluation module.
  Then we have
  $J^n \Phi(z)
    = \Phi(z) \Delta_{z, 0}(J^n)$, where
  \[
    \Delta_{z, 0}(J^n)
    = \oint_z \frac{d\xi}{2\pi i} \xi^n
    \bigg(
      \sum_m (\xi - z)^{-m - 1} y^m
    \bigg) \otimes 1 + 1 \otimes J^n.
  \]
  Note that $\oint_0 J(z) = \sum_n J_n z^{-n - 1}$.
  In vertex operator algebras,
  there is an expansion
  \[
    A(\xi) B(z)
    = \sum \frac{A_n B(z)}{(\xi - z)^{n + 1}},
  \]
  and these correspond to
  shifting certain vector spaces ``sitting at
  $z$'' like above.
\end{remark}

\section{Hopf Algebras}

\begin{remark}
  A good reference for quantum groups is
  \emph{A guide to quantum groups} by
  Chari and Pressley.
\end{remark}

\begin{definition}
  A \emph{Hopf algebra} an algebra $A$ with
  \begin{itemize}
    \item a product
      $m : A \otimes A \to A$,
    \item a coproduct
      $\Delta : A \to A \otimes A$,
    \item a unit $i : \C \to A$,
    \item a counit $\varepsilon : A \to \C$,
    \item an antipode $\gamma : A \to A$,
  \end{itemize}
  such that the following
  are satisfied:
  \begin{enumerate}
    \item properties of the product and
      coproduct:
      \begin{enumerate}[(a)]
        \item $m(x \otimes m(y \otimes z)) = m(m(x \otimes y) \otimes z)$
          (or more simply, $m({\id} \otimes m) = m(m \otimes \id)$);
        \item $(\Delta \otimes \id) \Delta = ({\id} \otimes \Delta) \Delta$;
      \end{enumerate}
    \item properties of the unit and counit:
      \begin{enumerate}[(a)]
        \item $m(i(1) \otimes x) = m(x \otimes i(1)) = x$;
        \item $(\varepsilon \otimes \id) \Delta = ({\id} \otimes \varepsilon) \Delta = \id$;
      \end{enumerate}
    \item properties of the antipode:
      $m(({\id} \otimes \gamma)(\Delta(x)))
      = m((\gamma \otimes \id)(\Delta(x)))
      = i(\varepsilon(x))$;
    \item consistency of product and coproduct:
      $\Delta(m(x \otimes y)) = m \otimes m(P_{2, 3}(\Delta(x) \otimes \Delta(y)))$, where
      $\Delta(a) = \sum a^{(1)} \otimes a^{(2)}$
      and
      \[
        (xy)^{(1)} \otimes (xy)^{(2)}
        = m(x^{(1)} \otimes y^{(1)}) \otimes m(x^{(2)} \otimes y^{(2)});
      \]
    \item consistency of unit and counit:
      $\varepsilon(i(1)) = 1$,
    \item consistency of product and
      counit (and of coproduct and unit):
      \begin{enumerate}[(a)]
        \item $\varepsilon(m(x \otimes y)) = \varepsilon(x) \varepsilon(y)$;
        \item $\Delta(i(1)) = i(1) \otimes i(1))$;
      \end{enumerate}
    \item consistency of product
      (and of coproduct) and antipode:
      \begin{enumerate}[(a)]
        \item $\gamma(m(x \otimes y)) = m^\mathrm{op}(\gamma(x) \otimes \gamma(y))$,
          where $m^\mathrm{op}(x \otimes y) = m(y \otimes x)$;
        \item $\Delta(\gamma(x)) = (\gamma \otimes \gamma)(\Delta^\mathrm{op}(x))$, where $\Delta^\mathrm{op}(x) = P_{1, 2}(\Delta(x))$
          and $P_{1, 2} : x \otimes y \mapsto y \otimes x$;
      \end{enumerate}
    \item consistency of unit (and of counit)
      and antipode:
      \begin{enumerate}[(a)]
        \item $\gamma(i(1)) = i(1)$;
        \item $\varepsilon(\gamma(x)) = \varepsilon(x)$.
      \end{enumerate}
  \end{enumerate}
\end{definition}

\begin{remark}
  We have the following:
  \begin{enumerate}
    \item Any Hopf algebra is an
      associative algebra with unit, and
      an associative coalgebra with counit.
    \item To describe the Hopf structure for
      an associative algebra with unit, it is
      enough to describe the action and relations
      on generators. Moreover,
      The unit, counit, and antipode
      are uniquely determined by the
      multiplication and comultiplication.
  \end{enumerate}
\end{remark}

\begin{corollary}
  We have the following:
  \begin{enumerate}
    \item If $A$ is a finite dimensional
      Hopf algebra, then $A^*$ is also a
      Hopf algebra with $m' = \Delta^*$, etc.
    \item Let $\Delta \to \Delta^{\mathrm{op}}$
      define another Hopf algebra. If
      $A$ is commutative, then $A^{\mathrm{op}} \cong A$
      with isomorphism $\gamma$.
    \item If $A, B$ are Hopf algebras,
      then $A \otimes B$ is also a Hopf algebra.
  \end{enumerate}
\end{corollary}

\begin{example}
  We have the following:
  \begin{enumerate}
    \item Let $G$ be a group and
      $A = \C[G]$. Then we can define
      a Hopf algebra structure on $A$ by
      \[
        m(g \otimes h) = gh, \quad
        \Delta(g) = g \otimes g, \quad
        \gamma(g) = g^{-1}, \quad
        \varepsilon(g) = 1, \quad
        i(1) = e,
      \]
      If $G$ is finite, we can also define
      the dual Hopf algebra $A^*$ to be
      $\mathcal{F}(G)$ (functions on $G$) with
      \[
        m(f \otimes \phi)(g)
        = f(g)\phi(g), \quad
        \Delta(f)(g, h) = f(g, h), \quad
        i(1) = 1, \quad
        \varepsilon(f) = f(e),
        \quad \gamma(f)(g) = f(g^{-1}).
      \]
    \item Let $G$ be a linear algebraic group.
      Then $A = \mathcal{P}(G)$ (the space of polynomial
      functions on $G$) is a Hopf algebra
      with the same formulas as above
      for $\mathcal{F}(G)$ in the finite
      group case.
    \item Let $\g$ be a Lie algebra. Then
      $U(\g)$ is a Hopf algebra with
      \[
        \Delta(x) = x \otimes 1 + 1 \otimes x, \quad
        \varepsilon(x) = 0, \quad
        \gamma(x) = -x.
      \]
  \end{enumerate}
  Note that the examples (2) and (3) are dual
  to each other.
\end{example}

\begin{definition}
  A \emph{representation} of a Hopf algebra
  $A$ is a left module over $A$ as an
  associative algebra.
\end{definition}

\begin{example}
  We can take $x \lambda = \varepsilon(x) \lambda$,
  where $\lambda$ is a $\C$-valued representation.
\end{example}

\begin{remark}
  Let $V$ be a finite-dimensional representation
  of $A$. Then there are two dual modules
  $V^*$ and ${}^*V$ for $A$ given by
  \begin{itemize}
    \item $A$ acts on $V^*$ by
      $a f(v) = f(\gamma(a) v)$ for
      $v \in V$, $f \in V^*$;
    \item $A$ acts on ${}^* V$ by
      $a f(v) = f(\gamma^{-1}(a) v)$.
  \end{itemize}
\end{remark}

\begin{exercise}
  Show that ${}^*V^*$ is canonically
  isomorphic to $V$, but ${}^{**}V$ and
  $V^{**}$ are different in general.
\end{exercise}

\begin{remark}
  Note that $V^* \otimes V \to \C$ and
  $\C \to V \otimes V^*$ are $A$-homomorphisms.
\end{remark}

\begin{remark}
  Let $\rho_i : A \to \End(V_i)$ for
  $i = 1, 2$. Define
  $\rho : A \to \End(V_1 \otimes V_2)$ by
  \[
    \rho(x) = \rho_1 \otimes \rho_2(\Delta(x)).
  \]
  Then $\Delta(\Delta \otimes I) = (I \otimes \Delta) \Delta$
  implies associativity of the tensor product.
\end{remark}

\section{Quantum Groups}
\begin{remark}
  Let $A$ be a generalized Cartan matrix
  with $d_i a_{i, j} = d_j a_{j, i}$. Fix
  $\{d_i\}$ to be positive integers
  with greatest common divisor $1$. Then
  we get an extended Lie algebra
  $\g_e$ with nondegenerate form
  $\llparenthesis \cdot, \cdot \rrparenthesis$.
  Normalize the form $\llparenthesis \cdot, \cdot \rrparenthesis$
  so that
  \[
    d_i = \frac{\llparenthesis \alpha_i, \alpha_i \rrparenthesis}{2},
  \]
  or equivalently, $\llparenthesis \alpha, \alpha \rrparenthesis = 2$ for short
  roots $\alpha$. Then
  $\mathfrak{h}_e \cong \mathfrak{h}^*_e$, and
  we can define
  \[
    h_i = \frac{2 \alpha_i}{\llparenthesis \alpha_i, \alpha_i \rrparenthesis}.
  \]
  Write $q = e^t$ for $t \in \C$ and
  $q^x = e^{tx}$ (assume that $q$ is not
  a root of unity). Then $U_q(\g)$ is generated
  by $e_i, f_i$ for $1 \le i \le r$ and
  $q^h$ for $h \in \mathfrak{h}$, with relations
  \begin{itemize}
    \item $q^0 = 1$ and
      $q^{a + b} = q^a q^b$
      for $a, b \in \mathfrak{h}$;
    \item $q^h e_i q^{-h} q^{\alpha_i(h)} e_i$
      and
      $q^h f_i q^{-h} = q^{-\alpha_i(h)} f_i$;
    \item $\displaystyle e_i f_j - f_j e_i = \delta_{i, j} \frac{q^{d_i h_i} - q^{-d_i h_i}}{q^{d_i} - q^{-d_i}}$;
    \item \emph{$q$-Serre relations}:
      $\displaystyle \sum_{n = 0}^{1 - a_{i, j}} \frac{(-1)^n}{[n]_i! [1 - a_{i, j} - n]_i!} e_i^n e_j e_i^{1 - a_{i, j} - n} = 0$,
      where
      \[
        [n]_i
        = \frac{q^{nd_i} - q^{-nd_i}}{q^{d_i} - q^{-d_i}}
        \quad \text{and} \quad
        [n]_i! = \prod_{p = 1}^n [p]_i,
      \]
      and identical relations hold for the $f_i$.
  \end{itemize}
  Note that we can formally write the
  $q$-Serre relations as
  $(\ad^{(q)}_{e_i})^{1 - a_{i, j}}(e_j) = 0$,
  where
  \[
    [x, y]_q
    = \ad_x^{(q)}(y)
    = xy - q^{\llparenthesis \alpha, \beta \rrparenthesis} yx.
  \]
\end{remark}
